% from ega1-5-fr.tex
\oldpage[I]{139}
préschéma $U\times_{\mathbf{Z}}V$
(\hyperref[I.3.2.7-fr]{3.2.7}), et est par suite un schéma affine
(\hyperref[I.3.2.2-fr]{3.2.2})~; comme $\Gamma_f(X)$ est fermé dans
$X\times_{\mathbf{Z}}Y$ (\hyperref[I.5.4.3-fr]{5.4.3}),
$\Gamma_f(X)\cap p_1^{-1}(U)\cap p_2^{-1}(V)$ est fermé dans
$U\times_{\mathbf{Z}}V$, et par suite le préschéma induit par le
sous-préschéma de $X\times_{\mathbf{Z}}Y$ associé à $\Gamma_f$
(\hyperref[I.4.2.1-fr]{4.2.1}), sur la partie ouverte
$\Gamma_f(X)\cap p_1^{-1}(U)\cap p_2^{-1}(V)$ de son espace sous-jacent, est
un sous-préschéma fermé d'un schéma affine, donc un schéma affine
(\hyperref[I.4.2.3-fr]{4.2.3}). La proposition résulte alors du fait que
$\Gamma_f$ est une immersion.

\begin{examples}[5.5.11]
\label{I.5.5.11-fr}
Le préschéma de l'exemple (\hyperref[I.2.3.2-fr]{2.3.2})
(<<~droite projective sur un corps $K$~>>) est \emph{séparé}, car pour le
recouvrement $(X_1,X_2)$ de $X$ par des ouverts affines,
$X_1\cap X_2=U_{12}$ est affine et $\Gamma(U_{12},\mathscr{O}_X)$, anneau des
fractions rationnelles de la forme $f(s)/s^m$ avec $f\in K[s]$, est engendré
par $K[s]$ et par $1/s$, donc les conditions de
(\hyperref[I.5.5.6-fr]{5.5.6}) sont vérifiées.

Avec le même choix de $X_1$, $X_2$, $U_{12}$ et $U_{21}$ que dans l'exemple
(\hyperref[I.2.3.2-fr]{2.3.2}), prenons cette fois pour $u_{12}$
l'isomorphisme qui à $f(s)$ fait correspondre $f(t)$~; on obtient cette fois
par recollement un préschéma \emph{intègre non séparé} $X$, car la première
condition de (\hyperref[I.5.5.6-fr]{5.5.6}) est vérifiée, mais non la seconde.
Il est immédiat ici que
$\Gamma(X,\mathscr{O}_X)\to\Gamma(X_1,\mathscr{O}_X)=K[s]$ est un
isomorphisme~; l'isomorphisme réciproque définit un morphisme
$f:X\to\operatorname{Spec}(K[s])$ qui est surjectif, et pour tout
$y\in\operatorname{Spec}(K[s])$ tel que $\mathfrak{j}_y\neq(0)$,
$f^{-1}(y)$ est réduit à un point, mais pour $\mathfrak{j}_y=(0)$,
$f^{-1}(y)$ se compose de \emph{deux} points distincts (on dit que $X$ est la
<<~droite affine sur $K$, où le point $0$ est dédoublé~>>).

On peut aussi donner des exemples où \emph{aucune} des deux conditions de
(\hyperref[I.5.5.6-fr]{5.5.6}) n'est vérifiée. Remarquons d'abord que dans le
spectre premier $Y$ de l'anneau de polynômes $A=K[s,t]$ à deux indéterminées
sur un corps $K$, l'ouvert $U$ réunion de $D(s)$ et de $D(t)$ \emph{n'est pas
un ouvert affine}. En effet, si $z$ est une section de $\mathscr{O}_Y$
au-dessus de $U$, il existe deux entiers $m\geq0$, $n\geq0$ tels que $s^m z$
et $t^n z$ soient les restrictions à $U$ de polynômes en $s$ et $t$
(\hyperref[I.1.4.1-fr]{1.4.1}), ce qui n'est évidemment possible que si la
section $z$ se prolonge en une section au-dessus de $Y$ tout entier,
identifiée à un polynôme en $s$ et $t$. Si $U$ était un ouvert affine, le
morphisme d'injection $U\to Y$ serait donc un isomorphisme
(\hyperref[I.1.7.3-fr]{1.7.3}), ce qui est absurde puisque $U\neq Y$.

Cela étant, prenons deux schémas affines $Y_1$, $Y_2$, spectres premiers des
anneaux $A_1=K[s_1,t_1]$, $A_2=K[s_2,t_2]$~; prenons
$U_{12}=D(s_1)\cup D(t_1)$, $U_{21}=D(s_2)\cup D(t_2)$, et pour $u_{12}$ la
restriction à $U_{21}$ de l'isomorphisme $Y_2\to Y_1$ correspondant à
l'isomorphisme d'anneaux qui à $f(s_1,t_1)$ fait correspondre $f(s_2,t_2)$~;
on a ainsi un exemple où aucune des conditions de
(\hyperref[I.5.5.6-fr]{5.5.6}) n'est satisfaite (le préschéma intègre ainsi
obtenu est dit <<~plan affine sur $K$, où le point $0$ est dédoublé~>>).
\end{examples}

\begin{remark}[5.5.12]
\label{I.5.5.12-fr}
\normalfont Étant donnée une propriété \textbf{P} de morphismes de préschémas,
considérons les propositions suivantes~:
\begin{enumerate}
  \item[\rm(i)] \emph{Toute immersion fermée possède la propriété \textbf{P}.}
  \item[\rm(ii)] \emph{Le composé de deux morphismes possédant la propriété
    \textbf{P} possède la propriété \textbf{P}.}
  \item[\rm(iii)] \emph{Si $f:X\to X'$, $g:Y\to Y'$ sont deux $S$-morphismes
    possédant la propriété \textbf{P}, $f\times_S g$ possède la propriété
    \textbf{P}.}

